% This template intentionally does NOT use the beamer class. We have to use ltx-talk 
% for accessibility

\DocumentMetadata{
	tagging=on,
	pdfstandard = ua-2,
	pdfversion  = 2.0,
	lang        = en-US,
	tagging-setup={
		math/alt/use,
		math/setup=mathml-SE
	}
}

\documentclass[aspect-ratio=16:9]{ltx-talk}

% ltx-talk tags every \frametitle at heading level H4 by default (the
% relationship between frame titles and document sectioning is still an
% open design question upstream -- see the manual, section 14.2), even
% though \section still correctly tags as H1. On a frame with no
% preceding \section (e.g. an early "Outline" slide), that H4 is the
% first heading in the document, which checkers such as Blackboard Ally
% flag as "headings do not begin at level 1." Promoting frametitle to H1
% removes that jump; \section headings stay H1 too, so a deck reads as a
% flat sequence of H1 headings rather than a nested outline, but that
% trades away a distinction checkers don't require in exchange for one
% they do.
\tagpdfsetup{role/new-tag=frametitle/H1}

\usepackage{tikz}
\usetikzlibrary{arrows.meta, positioning, calc, shapes.geometric}
% amsmath and its math fonts (via unicode-math for LuaTeX) are already
% loaded by ltx-talk; loading amssymb here would redefine symbols that
% unicode-math already provides and error out.
\usepackage{booktabs}
\usepackage{array}

% --- William & Mary color palette ---
% Hex values follow the live official palette at
% wm.edu/brand/visual-brand/color-palette (fetched 2026-07-21), which
% supersedes the older hex values in third-party templates such as
% https://github.com/leadmocha/beamer_template_WM . Per that page,
% Primary/Secondary colors should cover ~70-80% of the color space;
% Tertiary colors are accents only, capped at ~10-15%. Spirit Gold is
% never paired with W&M Gold in the same design.
\definecolor{WMGreen}{HTML}{004E38}      % Primary: core brand green
\definecolor{WMGold}{HTML}{B79257}       % Primary: core brand gold
\definecolor{WMSpiritGold}{HTML}{FDB714} % Spirit palette accent; not used with WMGold
\definecolor{WMSilver}{HTML}{D8DCDB}     % Primary: supporting neutral

% Secondary/Tertiary accents for the "better vs. weaker" feedback example
% later. GriffinGreen is an official Secondary green (an enhancement
% tone, distinct from WMGreen); BrickRed is an official Tertiary accent.
% Both are used here, rather than WMGold, because WMGold on white is
% only about 2.9:1 contrast and fails WCAG AA (needs 4.5:1) -- see the
% note below.
\definecolor{GriffinGreen}{HTML}{28463D}
\definecolor{BrickRed}{HTML}{964A37}

% ltx-talk exposes its role colors through \DeclareColor rather than
% beamer's \setbeamercolor. "structure" drives frame titles, block
% titles, itemize bullets, and table-of-contents entries; "alert" drives
% \alert{...}; "example" is a second accent used below for a contrasting
% block. WMGold is kept out of this list on purpose: at ~2.9:1 contrast
% against white it fails WCAG AA (needs 4.5:1) even set in bold, so it
% only appears here as a non-text accent (the TikZ diagram's warning
% color), never as the sole color of readable text.
\DeclareColor{structure}{WMGreen}
\DeclareColor{alert}{BrickRed}
\DeclareColor{example}{GriffinGreen}

% Start frame content right under the title instead of the class default
% of vertically centering it in the remaining space.
\SetKeys[talk]{frame/vertical-alignment=top}

% --- Logo -----------------------------------------------------------------
% Both marks were extracted from the official WM_Simple_PPT_Template_2026.potx
% (a .potx is just a zip archive; the brand assets live under ppt/media/) --
% see the README for the extraction command and the other marks that came
% out of it. The logo is purely decorative in both spots below -- the
% institution is already named as real text via \institute{} and the
% footer -- so it is tagged as a non-content artifact rather than a Figure
% needing alt text. Note the key is lowercase \texttt{tag=artifact}: the
% graphics-tagging code matches that literal string, and a capitalized
% \texttt{Artifact} silently falls through to being treated as a custom
% struct name that still demands alt text.
\newcommand{\wmLogoColorFile}{logo-assets/wm-logo-color.png}
\newcommand{\wmLogoHeight}{0.22cm}

% Small monogram tucked immediately left of the footer's date text on
% every content page (see \EditInstance{footer}{std} below). \Gm@lmargin
% (1cm here) is where that text starts, so the box below right-aligns the
% mark inside [0cm, 1cm - 0.08cm], leaving a small gap before the text.
% \IfFileExists keeps the deck compiling cleanly if the file is missing.
\newcommand{\wmPutLogoBox}[3]{% #1 file, #2 y, #3 box width measured from the left paper edge
	\IfFileExists{#1}
		{\put(0cm,#2){\makebox[#3][r]{\includegraphics[tag=artifact, height=\wmLogoHeight]{#1}}}}
		{}
}

% "current/total" page count, flush against the right margin. The
% built-in "framenumber" footer-element can't sit flush right: ltx-talk's
% footer template always appends one trailing \hfil after the last
% element (to let it distribute slack for centering), which pulls the
% last item back from the edge rather than pinning it there -- no
% element-order/separator combination avoids that. So this is drawn the
% same way as the corner logo above: an absolute box ending exactly at
% \Gm@rmargin (1cm), on the footer's own baseline. \@framenumber and
% \@totalframes are the same class-provided macros the (now unused)
% "framenumber"/"totalframes" footer-elements would have displayed.
\makeatletter
\newcommand{\wmPutPageCount}{%
	\put(0cm,-9.8cm){\makebox[\dimexpr\paperwidth-1cm\relax][r]{\color{WMGreen}\tiny\@framenumber/\@totalframes}}%
}
\makeatother

% A thin green top bar on every page, echoing the reference theme's plain
% frame background (a slim WMGreen rule along the top edge), plus the
% corner logo and page count above. This hooks into the kernel's
% shipout/background hook -- the same mechanism the class's own
% \pagecolor uses -- so it only paints the page canvas and never touches
% the tag tree. The footer's baseline sits 0.2cm above the bottom edge
% (bmargin 0.8cm minus footskip 0.6cm), hence -9.8cm below. Skipped on
% the title page (\ifwmTitlePage, toggled around \maketitle below): that
% page is filled solid WMGreen, and this bar/logo/count are drawn in
% WMGreen too, so they'd render but be genuinely invisible -- same color
% as the background underneath them.
\newif\ifwmTitlePage
\AddToHook{shipout/background}
	{\ifwmTitlePage\else
	 \color{WMGreen}\put(0cm,-3mm){\rule{\paperwidth}{3mm}}%
	 \wmPutLogoBox{\wmLogoColorFile}{-9.8cm}{0.92cm}%
	 \wmPutPageCount
	 \fi}

% Frame titles are set larger and bolder, closer to the reference's
% ~20pt frame-title scale. Color stays WMGreen rather than the
% reference's WMGold: gold frame titles on a plain white page only
% reach about 2.9:1 contrast, short of the 3:1 WCAG AA floor even for
% large bold text.
\EditInstance{frametitle}{header}{font=\bfseries\Large}

% A two-part footer (date, title) echoing the reference theme's footline
% layout, recolored WMGreen for visibility (the reference sets the
% page-number element to WMWhite on the plain page background, which
% would be unreadable outside its own colored footer band). The page
% count is no longer part of this templated footer -- see
% \wmPutPageCount above.
\EditInstance{footer}{std}
	{element-order={date,title}, color=WMGreen}

% \tableofcontents entries already carry the target slide number
% internally (see the .toc file: \contentsline{section}{Foundations}{3}{...}),
% but ltx-talk's \l@section leaves the two hook slots meant for
% rendering it -- contentsline/page/before and .../page/after -- empty
% by default, so no number ever prints. Filling in .../page/after with
% a dotted leader gives each entry a classic "Title .......... 3" line;
% #3 is the slide number (the hook's arguments are level, title, page,
% destination, per \__talk_toc_level:nnnn in the class source).
\AddToHookWithArguments{contentsline/page/after}{\dotfill\,#3}

% A two-tone banner, closer to the reference theme's gold-title /
% green-body block boxes, but keeping every text color at a safe
% contrast: black-on-gold (~7.6:1) and white-on-green (~8.5:1) both
% clear WCAG AA. Only the title becomes a colored banner; the body text
% below it is left as ordinary, plainly-tagged paragraph or list
% content rather than being wrapped in a colored box itself.
\newcommand{\wmblocktitle}[3]{%
	\par\noindent
	\colorbox{#1}{\parbox{\dimexpr\linewidth-2\fboxsep\relax}{\color{#2}\bfseries\large #3}}%
	\par\vspace{4pt}%
}

% White title-page text, since the title page now gets a full WMGreen
% background (see the \AddToHookNext call just before \maketitle).
\EditInstance{titlepage-element}{title}{color=white, font=\bfseries\Huge}
\EditInstance{titlepage-element}{subtitle}{color=white, font=\Large}
\EditInstance{titlepage-element}{author}{color=white}
\EditInstance{titlepage-element}{institute}{color=white}
\EditInstance{titlepage-element}{date}{color=white}

% A crest between the subtitle and author lines (see the element-order
% passed to \maketitle below), using the ornate crown-and-cipher mark
% extracted from the .potx rather than the flat monogram used in the
% footer -- see the README. ltx-talk's titlepage-element instances read
% their content from a metadata token list named after the element
% (\@title, \@subtitle, ...), so "logo" needs a matching \@logo.
\newcommand{\wmCipherFile}{logo-assets/wm-cipher-crest-white.png}
\DeclareInstance{titlepage-element}{logo}{talk}{before-skip=0.5em, after-skip=0.5em}
\makeatletter
\ExplSyntaxOn
\tl_new:N \@logo
\tl_gset:Nn \@logo
	{
		\IfFileExists{\wmCipherFile}
			{\includegraphics[tag=artifact, height=1.8cm]{\wmCipherFile}}
			{}
	}
\ExplSyntaxOff
\makeatother

\hypersetup{
	pdftitle={Accessible LuaLaTeX Presentation Template},
	pdfauthor={Author},
	pdfsubject={Example accessible ltx-talk presentation styled for William & Mary},
	pdfkeywords={accessibility, tagged PDF, PDF/UA-2, LuaLaTeX, ltx-talk, William & Mary}
}

\title{Accessible Presentations with LuaLaTeX}
\subtitle{A tagged-PDF Template}
\author{Author}
\institute{William \& Mary}
\date{Summer Semester}

\begin{document}

% Fill just the title page with WMGreen, echoing the reference theme's
% green title-page background. \AddToHookNext fires this exactly once,
% on the very next shipout, so later pages are unaffected. framestyle=
% plain drops the class's own std footer/header (date/title text,
% WMGreen) for this one page -- like the bar/logo/count above, it would
% otherwise render invisibly on top of the same-colored fill.
\wmTitlePagetrue
\AddToHookNext{shipout/background}
	{\color{WMGreen}\put(0cm,-\paperheight){\rule{\paperwidth}{\paperheight}}}
\maketitle[element-order={title,subtitle,logo,author,institute,date}, framestyle=plain]
\wmTitlePagefalse

\begin{frame}
	\frametitle{Outline}
	
	\vspace*{10mm}
	\tableofcontents
\end{frame}

\section{Foundations}

\begin{frame}
	\frametitle{Why This Template Exists}
	Slides often look polished while hiding their structure from assistive technology. This template compiles with LuaLaTeX to a tagged PDF/UA-2 file that exposes the same organization to a screen reader that a sighted viewer sees on screen.

	\wmblocktitle{WMGold}{black}{Key idea}
	An accessible slide deck exposes real headings, real lists, real tables, and real math, not just visual layout that looks like them.
\end{frame}

\begin{frame}
	\frametitle{A Note on Overlays}
	This template avoids \texttt{\textbackslash pause} and other overlay reveals on purpose.

	\begin{itemize}
		\item Each overlay step becomes a separate output page.
		\item A screen reader then encounters the same slide several times in a row, with only small differences.
		\item Every frame here is written as one complete, final page instead.
	\end{itemize}

	\wmblocktitle{BrickRed}{white}{If you do use overlays}
	Check the tagged reading order afterward: duplicated near-identical pages are confusing without sight.
\end{frame}

\section{Structure and Math}

\begin{frame}
	\frametitle{Headings Create Real Navigation}
	Every frame title in this deck is a real \texttt{\textbackslash frametitle}, and every \texttt{\textbackslash section} command both starts a PDF bookmark and groups related frames.

	\begin{itemize}
		\item Bookmarks let a reader jump between topics without scrolling.
		\item Tags let a screen reader announce a heading before a frame title, the same way it would for a printed article section.
		\item Neither happens automatically from font size or boldface alone.
	\end{itemize}

	\alert{Current limitation:} \texttt{ltx-talk} tags frame titles at heading level H4 rather than H1, since the relationship between frame titles and document sectioning is still being decided upstream. Some checkers may flag this.
\end{frame}

\begin{frame}
	\frametitle{Displayed Equations}
	Use math environments rather than raw display delimiters, so the tagging layer can attach a MathML-oriented alternative to each formula.

	\begin{equation}
		J(\theta)=\frac{1}{m}\sum_{i=1}^{m}\left(h_\theta(x_i)-y_i\right)^2
		\label{eq:objective}
	\end{equation}

	\begin{equation}
		\theta_{k+1}=\theta_k-\alpha\nabla J(\theta_k)
		\label{eq:gradient-update}
	\end{equation}
\end{frame}

\begin{frame}
	\frametitle{Lists Near Equations}
	A short list is often clearer than a dense paragraph when a formula has several symbols.

	\begin{itemize}
		\item $m$ is the number of observations.
		\item $h_\theta(x_i)$ is the model prediction for observation $i$.
		\item $y_i$ is the measured response.
		\item $\nabla J(\theta_k)$ is the gradient evaluated at iteration $k$.
	\end{itemize}
\end{frame}

\section{Tables and Figures}

\begin{frame}
	\frametitle{A Tagged Table}
	\begin{table}
		\centering
		\caption{Accessibility checks for a presentation PDF.}
		\label{tab:checklist}
		\tagpdfsetup{table/header-rows={1}}
		\begin{tabular}{>{\raggedright\arraybackslash}p{0.32\linewidth}
				>{\raggedright\arraybackslash}p{0.55\linewidth}}
			\toprule
			Check                 & Why it matters                                            \\
			\midrule
			Real frame titles     & Screen readers and PDF outlines rely on headings.          \\
			Figure alt text       & Visual information needs a text alternative.               \\
			Declared table header & Data cells need a header to remain meaningful out of order. \\
			Equation environments & Math needs semantic structure, not just glyphs.            \\
			\bottomrule
		\end{tabular}
		\tagpdfsetup{table/header-rows={}}
	\end{table}
\end{frame}

\begin{frame}
	\frametitle{Tagged Vector Figure: Workflow}
	The \texttt{alt} key on the \texttt{tikzpicture} environment gives the diagram a text alternative without changing how it looks.

	\begin{figure}
		\centering
		\begin{tikzpicture}[
			alt={Flow diagram showing a four step accessible presentation workflow:
					draft semantic source, compile with LuaLaTeX tagging, inspect reading
					order and alternative text, then publish the tagged PDF and source files.},
			node distance=1.35cm,
			font=\small
			]
			\tikzset{
				process/.style={draw=WMGreen, fill=WMGreen!8, rounded corners=2pt, very thick,
						align=center, text width=2.3cm, minimum height=1cm},
				check/.style={draw=GriffinGreen, fill=GriffinGreen!10, rounded corners=2pt,
						very thick, align=center, text width=2.2cm, minimum height=1cm},
				warn/.style={diamond, draw=WMGold, fill=WMGold!15, very thick,
						aspect=2, align=center, inner sep=1pt, text width=2.1cm},
				flowarrow/.style={-{Stealth[length=3mm]}, very thick, draw=black!60}
			}
			\node[process] (draft) {\small Draft semantic\\source};
			\node[process, right=of draft] (compile) {\small Compile with\\LuaLaTeX};
			\node[warn, right=of compile] (inspect) {Inspect\\output};
			\node[check, right=of inspect] (publish) {Publish PDF\\and source};

			\draw[flowarrow] (draft) -- node[above, font=\scriptsize] {structure} (compile);
			\draw[flowarrow] (compile) -- node[above, font=\scriptsize] {tags} (inspect);
			\draw[flowarrow] (inspect) -- node[above, font=\scriptsize] {fixes} (publish);
			\draw[flowarrow, bend right=-45] (inspect.south) to node[below, font=\scriptsize] {revise} (draft.south);
		\end{tikzpicture}
	\end{figure}
\end{frame}

\begin{frame}
	\frametitle{Writing Useful Alternative Text}
	\begin{columns}
		\begin{column}{0.48\textwidth}
			\wmblocktitle{GriffinGreen}{white}{Better}
			Contour plot of a bowl-shaped objective with a descent path moving toward the minimum.
		\end{column}
		\begin{column}{0.48\textwidth}
			\wmblocktitle{BrickRed}{white}{Weaker}
			Image of a plot with green and gold lines.
		\end{column}
	\end{columns}

	\vspace{0.4cm}
	Good alternative text names the figure type, states the main pattern or conclusion, and includes labels or trends the argument depends on. It rarely needs to list every visual detail.
\end{frame}

\section{Before You Share}

\begin{frame}
	\frametitle{Validation Checklist}
	\begin{columns}
		\begin{column}{0.48\textwidth}
			\wmblocktitle{WMGold}{black}{Before Compiling}
			\begin{itemize}
				\item Real frame titles and sections
				\item Equation environments, not raw \texttt{\$\$}
				\item Alt text for every meaningful figure
				\item Header row declared for every table
			\end{itemize}
		\end{column}
		\begin{column}{0.48\textwidth}
			\wmblocktitle{WMGold}{black}{After Compiling}
			\begin{itemize}
				\item Compiled with \texttt{lualatex}, twice
				\item PDF reports as tagged
				\item Title, author, and language are set
				\item Reading order matches the visual order
			\end{itemize}
		\end{column}
	\end{columns}
\end{frame}

\begin{frame}
	\frametitle{References}
	\begin{itemize}
		\item World Wide Web Consortium. \emph{Web Content Accessibility Guidelines (WCAG) 2.1}. W3C Recommendation, 2018.
		\item International Organization for Standardization. \emph{ISO 14289-2:2024, Document management applications -- Electronic document file format enhancement for accessibility -- Part 2: Use of ISO 32000-2 (PDF/UA-2)}. ISO, 2024.
	\end{itemize}
\end{frame}

\begin{frame}
	\frametitle{Questions?}
	
	\vspace*{10mm}
	\textbf{Source:} \texttt{https://github.com/dvasiliu/Accessibility_Examples}
\end{frame}

\end{document}
